\documentclass[a4paper,12pt]{article}
\usepackage[margin=1in]{geometry}
\usepackage{parskip}

\begin{document}

October 5, 2026

Transportation Research Part D: Transport and Environment



~

~

Dear Editors,

Please find our manuscript, titled ``Consumer Valuation of Battery Health Attributes in Used Electric Vehicles,'' submitted for consideration in \textit{Transportation Research Part D}.

Used battery electric vehicles (BEVs) are entering the secondary market at a growing pace, yet buyers face substantial uncertainty about battery condition at the point of sale. The battery can account for up to 30\% of a used BEV's value, and unlike vehicle age or mileage, battery health is not directly observable. This information gap depresses resale prices and may suppress new BEV adoption by raising concerns about future residual value. 

This study delivers the first estimate of consumers' willingness to pay (WTP) for battery-health attributes in the used BEV market, using a discrete choice experiment (DCE) administered to a national sample of 3,072 potential used vehicle buyers in the U.S.. The DCE jointly varies purchase price, electric range, vehicle mileage, battery degradation rate, and battery refurbishment history across six repeated choice tasks. In addition, respondents were randomly assigned to receive either basic battery information or extended information that additionally describes replacement costs and warranty coverage, allowing us to test whether disclosure shapes preferences and market engagement.

At the aggregate level, consumers value electric range, penalize mileage and battery degradation, and discount refurbishment history, but the preference heterogeneity is substantial. The LCCM results suggest that consumers fall into six segments that differ markedly in their sensitivity to battery quality signals, their price responsiveness, and their market engagement. Findings from this study show the potential for battery-health disclosure and certification to serve as a market-transparency signal in used BEV transactions. For dealers and OEMs, state of health, degradation history, and remaining warranty coverage could command measurable price premiums from the quality-sensitive majority.

This manuscript is original and unpublished. We look forward to receiving feedback.

\vspace{2em}

Sincerely,

\vspace{1em}

Xiatian Iogansen, Ph.D.\\
Department of Engineering Management and Systems Engineering, The George Washington University\\
Applied Economic Office, Engineering Laboratory, National Institute of Standards and Technology\\
Email: xiatian.iogansen@gwu.edu

\end{document}
